\section{Билет 24. Общее волновое уравнение. Скорость упругой волны в струне с выводом.}

\begin{definition}
\textbf{Общее волновое уравнение} для одномерного случая имеет вид:
\begin{equation}
    \frac{\partial^2 \xi}{\partial t^2} = v^2 \frac{\partial^2 \xi}{\partial x^2}, \label{eq:wave_1d_24}
\end{equation}
где $\xi(x,t)$ --- поперечное смещение элемента среды, $v$ --- фазовая скорость волны. Любая функция вида $f(x \pm vt)$ является решением этого уравнения и описывает бегущую волну.
\end{definition}

\begin{theorem}[Вывод скорости волны в струне]
Рассмотрим гибкую однородную струну с линейной плотностью $\rho$ (масса единицы длины), натянутую с силой $F$. Пусть в струне распространяется малая поперечная волна, смещение частиц $\xi(x,t)$ перпендикулярно оси $x$.

Выделим элемент струны длиной $dx$. Силы натяжения $F$ действуют на концах элемента под малыми углами $\alpha_1$ и $\alpha_2$ к оси $x$. Проекция результирующей силы на ось $\xi$:
\begin{equation}
    F_{\xi} = F \sin\alpha_2 - F \sin\alpha_1 \approx F(\tan\alpha_2 - \tan\alpha_1),
\end{equation}
так как при малых углах $\sin\alpha \approx \tan\alpha$. Но $\tan\alpha = \partial \xi / \partial x$ --- наклон касательной к струне. Следовательно:
\begin{equation}
    F_{\xi} = F \left[ \left.\frac{\partial \xi}{\partial x}\right|_{x+dx} - \left.\frac{\partial \xi}{\partial x}\right|_{x} \right] = F \frac{\partial^2 \xi}{\partial x^2} dx.
\end{equation}
Масса элемента $dm = \rho dx$. По второму закону Ньютона в проекции на ось $\xi$:
\begin{equation}
    (\rho dx) \frac{\partial^2 \xi}{\partial t^2} = F \frac{\partial^2 \xi}{\partial x^2} dx.
\end{equation}
Сокращая на $dx$, получаем волновое уравнение:
\begin{equation}
    \frac{\partial^2 \xi}{\partial t^2} = \frac{F}{\rho} \frac{\partial^2 \xi}{\partial x^2}.
\end{equation}
Сравнивая с \eqref{eq:wave_1d_24}, находим \textbf{скорость поперечной волны в струне}:
\begin{equation}
    v = \sqrt{\frac{F}{\rho}}. \label{eq:string_speed}
\end{equation}
\end{theorem}

\begin{property}
Из формулы \eqref{eq:string_speed} следуют важные практические выводы:
\begin{enumerate}
    \item Скорость волны \textbf{не зависит} от частоты и амплитуды колебаний (в пределах применимости линейного приближения).
    \item При увеличении силы натяжения $F$ скорость растёт: $v \propto \sqrt{F}$.
    \item При увеличении линейной плотности $\rho$ (более толстая или тяжёлая струна) скорость уменьшается: $v \propto 1/\sqrt{\rho}$.
\end{enumerate}
\end{property}

\begin{remark}
\textbf{Применение формулы:}
\begin{itemize}
    \item \textbf{Музыкальные инструменты:} Частота основной гармоники струны длиной $L$ с закреплёнными концами: $\nu_1 = v/(2L) = \frac{1}{2L}\sqrt{F/\rho}$. Настройка осуществляется изменением натяжения $F$.
    \item \textbf{Оценка параметров:} Для стальной струны ($\rho \approx 0.01$ кг/м) при натяжении $F = 100$ Н: $v \approx \sqrt{100/0.01} = 100$ м/с. При длине $L = 0.65$ м основная частота $\nu_1 \approx 100/(2 \cdot 0.65) \approx 77$ Гц.
    \item \textbf{Ограничения:} Формула справедлива при малых амплитудах, когда можно пренебречь изменением натяжения при колебаниях и считать струну идеально гибкой (не сопротивляющейся изгибу).
\end{itemize}
\end{remark}
